% ===================================================================
% 04_aplicacion.tex
% Sección 4: Aplicación del Módulo
% ===================================================================

\begin{frame}{Aplicación del Módulo - Herramientas y Tecnologías}
    \begin{columns}[t]
        \column{0.5\linewidth}
        \textbf{Herramientas Utilizadas}
        \begin{description}
            \item[Python 3.8] Lenguaje de programación
            \item[Pandas] ETL y manipulación de datos
            \item[Seaborn/Math] Visualización y análisis
            \item[scikit-learn] Modelado predictivo
            \item[Jupyter] Desarrollo interactivo
        \end{description}
        
        \column{0.5\linewidth}
        \begin{block}{Datos del Proyecto}
            \begin{itemize}
                \item Dataset: \datasetname
                \item Registros: 7043
                \item Variables tras procesamiento: 31
                \item Target: Churn (0 = No, 1 = Sí)
                \item Fuente: Kaggle
            \end{itemize}
        \end{block}
    \end{columns}
    
    \vsp
    \begin{exampleblock}{Variables Principales}
        \begin{columns}
            \column{0.33\linewidth}
            \textbf{Demográficas}
            \begin{itemize}
                \item Edad
                \item Género
                \item Estado civil
                \item Dependientes
            \end{itemize}
            \column{0.33\linewidth}
            \textbf{Servicios}
            \begin{itemize}
                \item Teléfono
                \item Internet
                \item TV
                \item Seguridad
            \end{itemize}
            \column{0.33\linewidth}
            \textbf{Financieras}
            \begin{itemize}
                \item Cargos mensuales
                \item Cargos totales
                \item Método de pago
                \item Contrato
            \end{itemize}
        \end{columns}
    \end{exampleblock}
\end{frame}

\begin{frame}{Aplicación del Módulo - ETL y Análisis Exploratorio}
    \begin{figure}
        \centering
        \begin{tikzpicture}[node distance=1.5cm, auto]
            % Etapas del proceso
            \node[rectangle, draw, fill=primary!20, minimum width=2.5cm, minimum height=0.8cm] (extrac) {Extracción};
            \node[rectangle, draw, fill=primary!20, right of=extrac, node distance=3cm, minimum width=2.5cm, minimum height=0.8cm] (clean) {Limpieza};
            \node[rectangle, draw, fill=primary!20, right of=clean, node distance=3cm, minimum width=2.5cm, minimum height=0.8cm] (transform) {Transformación};
            \node[rectangle, draw, fill=accent!20, right of=transform, node distance=3cm, minimum width=2.5cm, minimum height=0.8cm] (eda) {EDA};
            
            % Flechas
            \draw[->, thick, primary] (extrac) -- (clean);
            \draw[->, thick, primary] (clean) -- (transform);
            \draw[->, thick, accent] (transform) -- (eda);
            
            % Etiquetas
            \node[above of=extrac, node distance=0.5cm] {Datos crudos};
            \node[above of=clean, node distance=0.5cm] {Valores nulos};
            \node[above of=transform, node distance=0.5cm] {Codificación};
            \node[above of=eda, node distance=0.5cm] {Análisis};
        \end{tikzpicture}
        \caption{Flujo de procesamiento de datos}
    \end{figure}
    
    \begin{block}{Patrones Identificados}
        \begin{itemize}
            \item Contratos mes a mes → Mayor fuga
            \item Baja antigüedad → Mayor fuga
            \item Fibra óptica y cargos altos → Mayor propensión
            \item Ausencia de servicios adicionales → Mayor fuga
        \end{itemize}
    \end{block}
\end{frame}